The development of the sensor began with conductive TPU as the sensing material rather than the conductive PLA of the template. Conductive PLA was in fact the preferred starting point, since the template sensor on which the present design is based was realized and characterized in conductive PLA, so that its sensing behavior was already established and reusing the same material would have carried that behavior over with the least uncertainty. PLA is generally described as stiff and brittle, with a low glass-transition temperature and limited elongation at break~\cite{Farah2016,protopasta_conductive_pla}, and within additive-manufacturing practice it is commonly held that PLA becomes markedly more brittle below \SI{0}{\celsius}~\cite{PLAlowtemp}. This raised the concern that PLA would not function reliably toward the \SI{-20}{\celsius} end of the operating range. To avoid carrying that risk into the rest of the work, a flexible conductive TPU (NinjaTek EEL) was adopted first, an elastomer with a Shore hardness of 90A~\cite{ninjatek_ninjaflex_2026} whose elastomeric matrix was expected to better tolerate low temperatures. The low-temperature behavior of PLA is however not well documented, and this concern rested on widely repeated practical guidance rather than on measurements of the conductive PLA used here. The assumption, together with the choice between the two materials, is revisited in \Cref{Temperature Behaviour} once the temperature behavior has been measured.

Both materials were characterized by straining a single sensor and recording its resistance. The resistance at rest, before any strain is applied, is denoted the baseline resistance $R_0$, and the response to strain is expressed as the relative change $\Delta R / R_0$, which removes the dependence on the absolute baseline so that sensors with different baseline resistances can be compared on the same scale.

The sensor optimization process was developed in five stages, and the effect of each stage is shown for the conductive-TPU sensor in \cref{fig:opt_tpu} and for the conductive-PLA sensor in \cref{fig:opt_pla}, where the upper panel of each figure gives the resistance and the lower panel the relative change $\Delta R / R_0$, with one curve per stage.

The first optimization stage adjusts only the turn diameter and serves as the reference against which the later stages are compared. The second stage adds the line depth, which specifies how far the nozzle is lowered relative to the nominal sensing-layer height of \SI{0.2}{\milli\meter}, where a positive value brings the nozzle closer to the substrate and presses the deposited line more firmly while a negative value raises it, so that it sets the effective height and spread of the conductive line. This setting produces the largest single improvement of the sequence, and for the conductive PLA the turn-diameter-only stage barely responds, with a relative resistance change below \num{0.1}, which the line depth optimization raises to about \num{0.26}. The line depth therefore has a particularly strong effect on the sensitivity in the small-strain regime, and the relative resistance change increases further through the remaining stages for both materials, as summarized in \cref{tab:opt_results}.

\begin{table}
\centering
\caption{Resting-to-peak resistance range and relative resistance change $\Delta R/R_0$ of the conductive-TPU and conductive-PLA sensors after each print-optimization stage, evaluated at \SI{800}{\micro\strain}. Stages 3 and 4 each add a single turn-deposition setting, the speed or the flow, to the stage-2 configuration, and stage 5 applies both.}
\label{tab:opt_results}
\begin{tabular}{@{}c l c c c c@{}}
\toprule
 & & \multicolumn{2}{c}{TPU} & \multicolumn{2}{c}{PLA} \\
\cmidrule(lr){3-4}\cmidrule(lr){5-6}
Stage & Settings & $R$ [\si{\kilo\ohm}] & $\Delta R/R_0$ & $R$ [\si{\kilo\ohm}] & $\Delta R/R_0$ \\
\midrule
1 & Turn diameter  & 117--123 & 0.05 & 47--51 & 0.08 \\
2 & Stage 1 + line depth     & 130--146 & 0.12 & 39--49 & 0.26 \\
3 & Stage 2 + speed          & 132--152 & 0.15 & 37--48 & 0.31 \\
4 & Stage 2 + flow-rate      & 135--159 & 0.17 & 37--49 & 0.33 \\
%5 & Stage 1-4 & 112--135 & 0.20 & 36--50 & 0.38 \\
5 & Stage 2 + speed \& flow-rate & 112--135 & 0.20 & 36--50 & 0.38 \\
\bottomrule
\end{tabular}
\end{table}

\begin{figure}
\centering
\includestandalone[mode=buildmissing]{Standalone/Plots/TPU_optimized_sensors}
    \caption{Response of the conductive-TPU sensor after each of the five print-optimization stages, from the turn-diameter-only configuration to the combined speed and flow adjustment. The upper panel shows the measured resistance and the lower panel the relative change $\Delta R/R_0$.}
    \label{fig:opt_tpu}
\end{figure}

\begin{figure}
\centering
\includestandalone[mode=buildmissing]{Standalone/Plots/PLA_optimized_sensors}
    \caption{Response of the conductive-PLA sensor after the same five print-optimization stages as in \cref{fig:opt_tpu}.}
    \label{fig:opt_pla}
\end{figure}

% \begin{table}[!tb]
% \centering
% \caption{Resting-to-peak resistance range and relative resistance change $\Delta R/R_0$ of the conductive-TPU and conductive-PLA sensors after each print-optimization stage, evaluated at \SI{800}{\micro\strain}. Stages 3 and 4 each add a single turn-deposition setting, the speed or the flow, to the stage-2 configuration, and stage 5 applies both.}
% \label{tab:opt_results}
% \begin{tabular}{@{}c l c c c c@{}}
% \toprule
%  & & \multicolumn{2}{c}{TPU} & \multicolumn{2}{c}{PLA} \\
% \cmidrule(lr){3-4}\cmidrule(lr){5-6}
% Stage & Settings & $R$ [\si{\kilo\ohm}] & $\Delta R/R_0$ & $R$ [\si{\kilo\ohm}] & $\Delta R/R_0$ \\
% \midrule
% 1 & Turn diameter  & 117--123 & 0.05 & 47--51 & 0.08 \\
% 2 & Diameter+line depth     & 130--146 & 0.12 & 39--49 & 0.26 \\
% 3 & Diameter+line depth+speed          & 132--152 & 0.15 & 37--48 & 0.31 \\
% 4 & Diameter+line depth+flow-rate      & 135--159 & 0.17 & 37--49 & 0.33 \\
% 5 & All settings & 112--135 & 0.20 & 36--50 & 0.38 \\
% \bottomrule
% \end{tabular}
% \end{table}



% \begin{table}[!tb]
% \centering
% \small
% \caption{Optimized settings and the resulting resting-to-peak resistance range and relative resistance change $\Delta R/R_0$ of the conductive-TPU and conductive-PLA sensors after each print-optimization stage, evaluated at \SI{800}{\micro\strain}. A check mark marks each setting that is optimized at a given stage. Each stage keeps the settings of the earlier stages, except that stages~3 and~4 add the speed and the flow individually on top of stage~2, which stage~5 then combines.}
% \label{tab:opt_results}
% \begin{tabular}{@{}c cccc cc cc@{}}
% \toprule
%  & \multicolumn{4}{c}{Optimized setting} & \multicolumn{2}{c}{TPU} & \multicolumn{2}{c}{PLA} \\
% \cmidrule(lr){2-5}\cmidrule(lr){6-7}\cmidrule(lr){8-9}
% Stage & Diameter & Line depth & Speed & Flow & $R$ [\si{\kilo\ohm}] & $\Delta R/R_0$ & $R$ [\si{\kilo\ohm}] & $\Delta R/R_0$ \\
% \midrule
% 1 & \checkmark &            &            &            & 117--123 & 0.05 & 47--51 & 0.08 \\
% 2 & \checkmark & \checkmark &            &            & 130--146 & 0.12 & 39--49 & 0.26 \\
% 3 & \checkmark & \checkmark & \checkmark &            & 132--152 & 0.15 & 37--48 & 0.31 \\
% 4 & \checkmark & \checkmark &            & \checkmark & 135--159 & 0.17 & 37--49 & 0.33 \\
% 5 & \checkmark & \checkmark & \checkmark & \checkmark & 112--135 & 0.20 & 36--50 & 0.38 \\
% \bottomrule
% \end{tabular}
% \end{table}

The third and fourth stages both act on the deposition through the turns, where the over-extrusion that forms the contacts between adjacent turns most directly governs the sensitivity. Increasing the print speed moves the nozzle more quickly through the turn, so that less filament is deposited and the over-extrusion is reduced, while decreasing the flow rate has the same effect by reducing the amount of filament extruded. Applied on top of the turn diameter and line depth, the speed and flow-rate changes each add a smaller further improvement, and combining both in the fifth stage gives the best response, reaching a relative resistance change of about \num{0.38} for the conductive PLA. This configuration was adopted as the final print configuration for each material.
 
The same five stages were applied to both materials, with the settings optimized for the conductive TPU carried over to the conductive PLA as a starting template and then re-optimized for that material. The two sensors behave similarly through the stages, both showing the line depth as the dominant improvement, but the conductive TPU sits at a markedly higher baseline resistance and reaches only about half the relative resistance change of the conductive PLA, around \num{0.20} against \num{0.38} at full strain. The conductive TPU also exhibited more severe over-extrusion at the points between the turns than the PLA when no speed or flow compensation was applied. The speed and flow-rate changes required to correct the TPU were therefore larger than for the PLA. 
The final optimized print parameters for both materials, together with the finished sensing-element footprint, comprising the serpentine pattern and the contact pads used for wiring, are listed in \cref{tab:opt_params}.

\begin{table}
\centering
\caption{Final optimized print parameters for the conductive-TPU and conductive-PLA sensors.}
\label{tab:opt_params}
\begin{tabular}{@{}l c c@{}}
\toprule
Parameter & TPU & PLA \\
\midrule
Turn diameter [\si{\micro\meter}]            & 510 & 490 \\
Line depth [\si{\micro\meter}]               & 35 & 45 \\
Print speed in turns [\si{\milli\meter\per\second}] & 25 & 20 \\
Flow in turns [\si{\percent}]                & 90 & 98 \\
\midrule
Sensing-element footprint [\si{\milli\meter}]        & \multicolumn{2}{c}{$25 \times 6.25$} \\
\bottomrule
\end{tabular}
\end{table}

One further adjustment, Linear Advance, was attempted but is not shown in the figures because it did not improve the response. Linear Advance is a firmware feature reported by its developer to improve flow consistency during changes in print speed by compensating for the pressure that builds up in the nozzle~\cite{Lin_adv}. While its recommended calibration appeared promising, as shown in \cref{app:LinearAdvance}, applying it to the custom-generated serpentine G-code did not improve the flow but instead made it worse, most likely because the feature is intended for conventionally sliced G-code rather than the custom toolpaths used here.

